\chapter{Public Relations}
\label{ch:04_public-relations}

\begin{universityTask}
    Sie sind Leiter der Public-Relations-Abteilung einer politischen Partei und wollen als \q{Propagandamaßnahme} einen Social Bot einsetzen.

    \vspace{0.5cm}

    \begin{enumerate}[label=(\alph*)]
        \item Sind die auf diese Weise durchgeführten \q{Werbemaßnahmen} \q{above the line} oder \q{below the line}? Begründen Sie Ihre Antwort.
            \universityPoints{2}
        \item Nehmen wir an, das Hauptziel Ihrer Social-Bot-Propagandamaßnahme soll sein, das Corporate Image Ihrer Partei zu verbessern. Was versteht man überhaupt unter einem Corporate Image und wodurch wird dieses \q{Bild} wesentlich geprägt?
            \universityPoints{3}
        \item Nach welchen Kriterien könnten Sie die Zielgruppe für Ihre Propagandamaßnahme festlegen?
            \universityPoints{2}
        \item Skizzieren Sie grob eine Werbeleitstrategie, die der Social-Bot-Maßnahme zur Verbesserung des Corporate Image Ihrer Partei zugrunde liegen könnte.
            \universityPoints{5}
        \item Beschreiben Sie beispielhaft, wie die sechs Gestaltungstypen gemäß der systematischen Inhaltsanalyse bei Ihrer Propagandamaßnahme mit dem Social Bot eine Rolle spielen könnten.
            \universityPoints{6}
        \item Erläutern Sie den Zusammenhang zwischen der Social-Bot-Maßnahme und dem Begriff des \q{viralen Marketings}.
            \universityPoints{2}
    \end{enumerate}
\end{universityTask}

\section{Werbemaßnahmen}
\label{sec:04-01_advertising-measures}

Klassische Werbung ist in der Regel \Gls{atl} -- direkt erkennbar, oft unfreiwillig konsumiert und zielt auf eine breite Zielgruppe ab. Bei \Gls{btl} handelt es sich um Direktmarketing, das beispielsweise im Internet stattfindet. Mittels \Gls{btl} können spezifische und klein segmentierte Zielgruppen durch direkte Kommunikation erreicht werden. \cite[S. 26]{002_MKO03_communication-communication-goals-and-advertising-conception}

Ein Social Bot erfüllt mindestens zwei Merkmale von Direktmarketing:

\begin{itemize}
    \item Die Kommunikation erfolgt über das Medium Internet.
    \item Es findet eine direkte, adressaten- und zielgruppengerechte Kommunikation statt.
\end{itemize}

Der Social Bot ist in der Lage, durch Kommunikation mit dem Kunden für dessen individuelles Problem eine ebenso individuelle Produkt- oder Dienstleistungslösung zu idenifizieren. Social Bots ermöglichen theoretisch sowohl die direkte Kommunikation (mittels Direktnachrichten) als auch die indirekte Kommunikation (mittels Beiträgen oder Kommentaren) mit Kunden.

\section{Corporate Image}
\label{sec:04-02_corporate-image}

Das Corporate Image ist Teil der Corporate Identity und beschreibt das aktuelle Bild eines Unternehmens in den Köpfen der Menschen. Es wird wesentlich geprägt durch die Organisationsentwicklung, Personalentwicklung und Kommunikationspolitik eines Unternehmens. Dazu zählen sowohl Gestaltungs- und Werbemaßnahmen als auch Öffentlichkeitsarbeit. Unternehmensinterne Kommunikation trägt ebenfalls zum Corporate Image bei. \cite[S. 13 f.]{002_MKO03_communication-communication-goals-and-advertising-conception} Vereinfacht ausgedrückt schafft das Corporate Image also eine Vorstellung des Unternehmens bei den Menschen.

Für eine politische Partei könnte das Corporate Image beispielsweise durch die Farbe, das Logo und den Schriftzug geprägt sein, um den Wiedererkennungswert bei Bürgerinnen und Bürgern zu erhöhen. Darüber hinaus spielen die politische Position und inhaltliche Standpunkte sowie die Kontaktpunkte mit Bürgerinnen und Bürgern eine wichtige Rolle. Das Corporate Image einer Partei ist somit das in den Köpfen der Zielgruppe verankerte Fremdbild der Partei.

\section{Zielgruppe für Propagandamaßnahme}
\label{sec:04-03_target-group-for-propaganda-measure}

Die Zielgruppe definiert, an wen die Kommunikations- oder Werbemaßnahme adressiert ist. Die Beworbenen sind also bestenfalls auch die Zielgruppe. \cite[S. 6 f.]{002_MKO03_communication-communication-goals-and-advertising-conception}

Zielgruppen werden anhand verschiedener Kriterien unterschieden \cite[S. 20 ff.]{002_MKO03_communication-communication-goals-and-advertising-conception}:

\begin{itemize}
    \item Geschlecht, Alter und Einkommen sowie ähnliche Kennzahlen ordnen sich unter den \textbf{soziodemografischen Faktoren} ein.
    \item Die \textbf{geographischen Faktoren} beziehen sich beispielsweise auf die Großstadt, das Bundesland, die Bundesrepublik oder ordnen potentielle Kunden einem Kontinent beziehungsweise einer Region wie Europa zu.
    \item Ob ein Kunde beispielsweise interessiert an Kapitalmärkten ist, weil er die Zeitschrift \textit{DER AKTIONÄR} verfolgt, wäre ein Beispiel für \textbf{mediabezogene Faktoren}.
\end{itemize}


Neben diesen Kriterien können Zielgruppen anhand weiterer Eigenschaften voneinander unterschieden werden \cite[S. 23]{002_MKO03_communication-communication-goals-and-advertising-conception}:

\begin{itemize}
    \item \textbf{verhaltensdisponierend} sind
    \begin{itemize}
        \item \textbf{demografische} Merkmale (sozioökonomische und geographische)
        \item \textbf{psychografische} Merkmale (Motive, Einstellungen und Involvement)
    \end{itemize}
    \item \textbf{verhaltensdeskriptiv} ist das beobachtete \textbf{Kauf- und Konsumverhalten}, also
    \begin{itemize}
        \item \textbf{Kaufverhalten},
        \item \textbf{Konsumverhalten}, und
        \item \textbf{Besitzmerkmale}.
    \end{itemize}
\end{itemize}

Da es bei \textbf{Propagandamaßnahmen} darauf ankommt, die Zielgruppe möglichst in ihrem alltäglichen und gesellschaftlichen Leben zu erreichen, sind insbesondere \textbf{verhaltensdisponierende Maßnahmen} für die Zielgruppenanalyse von Relevanz. Sozioökonomische Merkmale wie Einkommen oder Bildungsgrad ermöglichen eine adressatengerechte Ansprache während geographische Merkmale für wahlkreisbezogene Einflüsse oder Themen relevant sein können. Wie genau kommuniziert werden sollte, lässt sich anhand von Werthaltungen und Überzeugungen sowie deren jeweiliger Ausprägung innerhalb der Zielgruppe ableiten.

\section{Werbeleitstrategie}
\label{sec:04-04_advertising-strategy}

Die Vorbereitung einer Werbeleitstrategie kann mittels einer sogenannten Copy-Strategie, auch Kreativbriefing genannt, erfolgen. Dabei werden zentrale Fragen zum \textbf{\nameref{sec:04-04-01_creative-objective}} (also dem Kommunikationsziel), der \textbf{\nameref{sec:04-04-02_creative-strategy}}, dem \textbf{\nameref{sec:04-04-03_consumer-benefit}} (das Produktversprechen), der \textbf{\nameref{sec:04-04-04_reason-why}} (dem Grund für das Produktversprechen) und der \textbf{\nameref{sec:04-04-05_tonality}} (der Stimmung oder Atmosphäre, auch Flair genannt) beantwortet. \cite[S. 31 f.]{002_MKO03_communication-communication-goals-and-advertising-conception}

Für die Werbeleitstrategie der Propagandamaßnahme gemäß Aufgabenstellung sei einmal angenommen, dass Propagandamaßnahmen naturgemäß eher ein Werkzeug extremer Parteien sind. Die konkrete politische Ausrichtung einer solchen Partei soll für dieses Beispiel unerwähnt bleiben und keine Rolle spielen. Die entwickelte Werbeleitstrategie ist als plakativ und exemplarisch angedacht, um die Möglichkeiten einer solchen Strategie zu verdeutlichen.

\begin{table}[htp]
    \centering
    \begin{university-table}{width=\textwidth, colspec={l X[1,1]}}
        \textbf{} & \textbf{Wichtigste Aspekte} \\
        \textbf{Creative Objective} & 
            Aufbau von Reichweite, Aufmerksamkeit, Image-Stabilisierung und wahrgenommener Bürgernähe durch kontinuierliche, dialogorientierte Social-Media-Kommunikation. \\
        \textbf{Creative Strategy} & 
            Profilierung der Partei als konsequente, meinungsstarke und nahbare politische Alternative. \\
        \textbf{Consumer Benefit} & 
            \qit{Wir sagen die Wahrheit, vertreten Sie und handeln mutig.} \linebreak
            Einzigartigkeit in Mut, Konsequenz und \qit{wahrer} Vertretung der Bürgerinteressen. \\
        \textbf{Reason Why} & 
            Ohne dieses digitale Netzwerk gäbe es keine \qit{echte} Gegenöffentlichkeit, die sich gegen die etablierten Mächte stemmt. \\
        \textbf{Tonality} & 
            Emotionen-geladen, alltagsnah, polemisch, aber scheinbar vernünftig und \qit{bürgerlich}. \\
    \end{university-table}
    \caption{Überblick über die Werbeleitstrategie einer Propagandamaßnahme}
    \label{tab:04-04_overview-over-advertising-strategy-for-propaganda-measure}
\end{table}

Die wichtigsten Aspekte aus \autoref{tab:04-04_overview-over-advertising-strategy-for-propaganda-measure} werden im Folgenden ausführlich erläutert.

\subsection{Creative Objective}
\label{sec:04-04-01_creative-objective}

Der Social Bot soll das Image einer politischen Partei als \q{hart, aber notwendig} etablieren und gleichzeitig eine scheinbar bürgerfreundliche, nahbare und vertrauenswürdige Marke nach außen transportieren. Ziel ist es, durch permanente, scheinbar organische Präsenz in sozialen Netzwerken Sympathie, Identifikation und Loyalität zu generieren und kritische Informationen zu relativieren oder zu übertönen.

\subsection{Creative Strategy}
\label{sec:04-04-02_creative-strategy}

Die Partei wird als \q{starke Stimme des Volkes} dargestellt, die klar, entschieden und konsequent handelt -- und nicht von \q{Mainstream-Medien} verfälscht wird. Der Social Bot agiert dabei nicht als offener Bot, sondern als scheinbar authentischer Nutzerprofil-Account, der eigenständig beiträgt, indem er

\begin{itemize}
    \item seinen eigenen Status und seine eigenen Posts erstellt,
    \item unter Beiträgen kommentiert,
    \item auf Kritik reagiert, und
    \item andere Profile teilt oder teilen lässt.
\end{itemize}

Der Eindruck soll entstehen, dass es sich um viele \q{normale} Bürger oder aber Mitglieder der Partei handelt, die dieselbe politische Haltung teilen -- nicht um eine steuerbares Propaganda-Werkzeug.

\subsection{Consumer Benefit}
\label{sec:04-04-03_consumer-benefit}

Ein Produktversprechen könnte sein: \qit{Wir sind die einzige Partei, die die Wahrheit sagt, für Sie einsteht und nicht mit den anderen Parteien gemeinsame Sache macht.}

Für den \textbf{\Gls{usp}} gilt:

\begin{itemize}
    \item Partei scheint einzige Kraft, die Mut und Entschlossenheit zeigt.
    \item Zusätzlich werden andere Parteien als kontrastlose Masse dargestellt.
\end{itemize}

Für die \textbf{\Gls{uap}} gilt:

\begin{itemize}
    \item Der Social Bot verstärkt das unter \Gls{usp} beschriebene Bild.
    \item Zentrale Themen werden immer kontinuierlich wiederholt.
    \item Alle Themen werden verpackt in Kommentaren oder Posts.
\end{itemize}

\subsection{Reason Why}
\label{sec:04-04-04_reason-why}

Die Partei kann diese Strategie ideologisch damit begründen, dass

\begin{itemize}
    \item die \q{Wahrheit} von der etablierten Medienlandschaft verschleiert werde, und
    \item eine \q{starke Stimme} nötig sei, um die Interessen \q{des kleinen Mannes} durchzusetzen.
\end{itemize}

Der Social Bot wird als selbstorganisierte, digitale Bürger-Bewegung inszeniert. Es wirkt, als ob viele Menschen ganz von sich aus dieselbe Meinung teilen. Tatsächlich werden aber Inhalte, Themen und Feindbilder systematisch von der Partei vorgegeben und wiederholt.

\subsection{Tonality}
\label{sec:04-04-05_tonality}

Die Stimmung und Atmosphäre der Propagandamaßnahme ließe sich durch die folgenden Eigenschaften charakterisieren:

\begin{itemize}
    \item \textbf{Emotional aufgeladen}: Es werden häufig indirekt oder offen Angst, Wut und Empörung ausgedrückt (\acrshort{zB} \qit{So kann das nicht weitergehen!}).
    \item \textbf{Scheinbar spontan und persönlich}: Es werden kurze, alltagsnahe Formulierungen und lokalisierende Beispiele verwendet (\acrshort{zB} \qit{in meinem Stadtteil ...}, \qit{in meiner Straße ...}).
\end{itemize}

Die Tonality ist differenziert genug, um als \q{vernünftiger Bürger} zu wirken, aber mit klaren, wiederkehrenden Feindbildern (\acrshort{zB} \qit{die da oben}, \qit{die Mainstream-Medien}, \qit{die etablierten Parteien}). Sobald Kritik auftritt, wird die Polemik aggressiver, aber nie so extrem, dass sie sofort als \q{Hass} oder \q{Hetze} erkennbar werden könnte.

\section{Gestaltungstypen der systematischen Inhaltsanalyse}
\label{sec:04-05_design-types-of-systematic-content-analysis}

Systematische Inhaltsanalyse, auch Content-Analyse genannt, beschäftigt sich mit der \qit{Erforschung des Zusammenhangs zwischen inhaltlichen Aussagen von Werbemitteln und deren Wirkung} \cite[S. 39 f.]{002_MKO03_communication-communication-goals-and-advertising-conception}. Es werden die folgenden sechs Faktoren unterschieden \cite[S. 39 f.]{002_MKO03_communication-communication-goals-and-advertising-conception}:

\begin{itemize}
    \item \textbf{\nameref{sec:04-05-01_suggestion}}: \linebreak Autoritäre Appelle bilden die Meinung entscheidungsschwacher Menschen.
    \item \textbf{\nameref{sec:04-05-02_identification}}: \linebreak Scheinbare Vorbilder überzeugen Menschen und wirken meinungsbildend.
    \item \textbf{\nameref{sec:04-05-03_presentation}}: \linebreak Aufwändig produzierte Werbemittel erregen Aufmerksamkeit.
    \item \textbf{\nameref{sec:04-05-04_concretion}}: \linebreak Realistische Darstellungen überzeugen Menschen mit geringer Abstraktionsfähigkeit.
    \item \textbf{\nameref{sec:04-05-05_information}}: \linebreak Informierende Werbung spricht Multiplikatoren an.
    \item \textbf{\nameref{sec:04-05-06_motivation}}: \linebreak Verlockende Versprechen sprechen den Status- oder Prestigewunsch von Menschen an.
\end{itemize}

Im Folgenden wird der Social Bot zur Propagandamaßnahme einer politischen Partei anhand dieser sechs Faktoren analysiert. Dabei wird auf der exemplarischen Werbeleitstrategie aus \autoref{sec:04-04_advertising-strategy} aufgebaut, um explorativ zu skizzieren, wie sich die Gestaltungstypen in der Praxis einer solchen Propagandamaßnahme zeigen könnten.

\subsection{Suggestion}
\label{sec:04-05-01_suggestion}

\begin{itemize}
    \item Druckvolle Posts wie \qit{Jetzt handeln, bevor es zu spät ist} oder \qit{Nur diese Partei schützt unsere Zukunft} regen zu sofortigen Reaktionen an.
    \item Sätze wie \qit{Es gibt nur noch eine echte Alternative} wirken absolut und lassen wenig Widerspruch zu.
    \item Emotional zugespitzte Aussagen über \qit{Bedrohung}, \qit{Verfall} oder \qit{Notlage} sollen Unsicherheit in klare Zustimmung umwandeln.
\end{itemize}

\subsection{Identifikation}
\label{sec:04-05-02_identification}

\begin{itemize}
    \item Geschickt eingebundene Zitate oder Geschichten von vermeintlichen Vorbildern wie Prominenten, Unternehmern oder bekannten Lokalfiguren wirken überzeugend.
    \item Beiträge in der Ich-Form wie \qit{Ich habe selbst erlebt, was in meinem Viertel passiert} schaffen Nähe und Wiedererkennung.
    \item Fotos oder kurze Videos mit scheinbar normalen Menschen aus ähnlichen Lebenslagen vermitteln das Gefühl: \qit{Das könnte auch ich sein}.
\end{itemize}

\subsection{Präsentation}
\label{sec:04-05-03_presentation}

\begin{itemize}
    \item Professionell geschnittene Kurzvideos, auffällige Grafiken oder plattformgerechte Memes ziehen Aufmerksamkeit auf sich.
    \item Emotional aufgeladene Bildstrecken mit starken Farben, Symbolen und Großaufnahmen verstärken die Wirkung.
    \item Regelmäßig veröffentlichte Posts mit hoher visueller Wiedererkennbarkeit lassen die Botschaft präsenter erscheinen.
\end{itemize}

\subsection{Konkretion}
\label{sec:04-05-04_concretion}

\begin{itemize}
    \item Konkrete Beispiele wie \qit{gestiegene Mieten im eigenen Stadtteil} oder \qit{überfüllte Busse morgens um sieben} machen politische Aussagen greifbar.
    \item Scheinbar persönliche Berichte über einen bestimmten Markt, eine Schule oder eine Straße wirken anschaulich und glaubwürdig.
    \item Einzelne Fälle von Problemen oder Konflikten werden so erzählt, als seien sie typisch für den Alltag vieler Menschen.
\end{itemize}

\subsection{Information}
\label{sec:04-05-05_information}

\begin{itemize}
    \item Posts mit Zahlen, kurzen Statistiken oder grafisch dargestellten Vergleichswerten vermitteln einen sachlichen Eindruck.
    \item Scheinbar neutrale Erklärbeiträge zu Themen wie Sicherheit, Preise oder Migration sprechen auch Personen an, die auf Fakten reagieren.
    \item Links zu Artikeln, Diagrammen oder Infokacheln können den Eindruck von Seriosität und Aufklärung verstärken.
\end{itemize}

\subsection{Motivation}
\label{sec:04-05-06_motivation}

\begin{itemize}
    \item Aussagen wie \qit{Wer sich engagiert, gehört zu den Mutigen} sprechen den Wunsch nach Anerkennung und Stärke an.
    \item Versprechen von Einfluss, Gemeinschaft oder \qit{Mitgestalten statt Zuschauen} aktivieren das Bedürfnis nach Zugehörigkeit.
    \item Formulierungen wie \qit{Sei Teil der Bewegung} oder \qit{Steh auf für das, was richtig ist} sollen zum Mitmachen und Positionieren anregen.
\end{itemize}

\section{Virales Marketing}
\label{sec:04-06_viral-marketing}

Virales Marketing nutzt existierende soziale Netzwerke und Medien, um Aufmerksamkeit auf Marken, Produkte oder Kampagnen zu lenken. Die Verbreitung der Nachrichten basiert dabei auf einer (elektronischen) Form der Mundpropaganda und funktioniert insbesondere im Medium Internet. \cite[S. 67 ff.]{002_MKO03_communication-communication-goals-and-advertising-conception}

Die Social-Bot-Maßnahme steht in engem Zusammenhang mit viralem Marketing, weil sie darauf abzielt, Inhalte schnell, sichtbar und möglichst breit im Netz zu verbreiten. Der Bot kann Posts automatisiert verbreiten, Kommentare setzen, Reaktionen auslösen und so den Eindruck erzeugen, eine Botschaft sei bereits besonders populär oder gesellschaftlich relevant. Dadurch wird Reichweite nicht nur gesteigert, sondern auch der Anschein organischer Zustimmung geschaffen.
	
Virales Marketing funktioniert grundsätzlich nach dem Prinzip, dass Nutzer Inhalte freiwillig weitergeben, weil sie diese interessant, emotional oder nützlich finden. Ein Social Bot unterscheidet sich davon, weil er diese Weiterverbreitung nicht spontan, sondern gezielt und technisch gesteuert unterstützt. Er ist damit kein klassischer Träger viralen Marketings, sondern ein Mittel, um virale Effekte künstlich anzuschieben, zu verstärken oder vorzutäuschen.
	
Ein kleines Beispiel wäre ein politischer Beitrag, der zunächst nur wenig Beachtung findet, aber durch viele Bot-Accounts in kurzer Zeit geliked, kommentiert und geteilt wird. Für Außenstehende wirkt der Beitrag dadurch wie ein bereits stark diskutiertes Thema, obwohl die Resonanz nicht unbedingt echt oder organisch entstanden ist.